\documentclass[11pt]{article}
\usepackage{fontspec}
\usepackage{amsmath,amssymb,booktabs,geometry,array,hyperref}
\geometry{a4paper,margin=18mm}
\setmainfont{WenQuanYi Micro Hei}
\hypersetup{hidelinks}
\setlength{\parskip}{0.45em}
\setlength{\parindent}{0pt}

\title{条件扩散的反事实归因\\[0.4em]
\large 一次只改一个坐标，分布差了多少}
\author{}
\date{2026年10月27日（太平洋时间）}

\begin{document}
\maketitle

一次归因问的是：其余条件保持这个窗口的观测值，只把条件向量里的一格换成参考路径，结果的条件分布离开当天观测多少。生成器是条件扩散模型。它在给定条件 \(c\) 时抽出一批结果，当作 \(\widehat P(Y\mid c)\)。分数是分布之间的距离比。它不是结构因果效应，也不把风控、推荐或社交通道拆开。

本文只写已经记下来的判定。扩散步数、噪声日程和网络宽度不在这份记录里，这里不补。北京东四看古城的查询在协议处停止，那一格没有训练扩散模型。

\section{编辑与分数}

条件向量在线性十站点验证面板上是 21 维，里面有协变量、干扰块，以及被怀疑的那一块滞后。对滞后 \(\ell\)，反事实条件 \(c^{(\ell)}\) 只替换这一格，其余坐标保持观测值。

观测批是这个窗口的真实结果。反事实批是模型在 \(c^{(\ell)}\) 上的抽样。无干预批是模型在原条件 \(c\) 上的抽样。同一评估日的归一化分数是
\begin{equation}
\label{eq:rho}
\rho(\ell)
=
\frac{\mathrm{MMD}^{2}\!\left(P^{\mathrm{obs}},\,\widehat P(Y\mid c^{(\ell)})\right)}
{\mathrm{MMD}^{2}\!\left(P^{\mathrm{obs}},\,\widehat P(Y\mid c)\right)}.
\end{equation}
分母是这一天不做这次编辑时，预测分布和观测已经有的距离。分子是改掉这一格之后的距离。核宽用训练期结果的边际中位距离，评估时不再另选。\(\rho(\ell)>2\) 是跑实验之前写死的线。

\(\rho\) 不带方向。头条句子用的是两批均值的符号差：正剂量要求观测批高于反事实批。距离 \(\rho\)、\(\Delta^{\mathrm{cf}}\) 和 \(\delta^{2}\) 在符号反过来的时候可以几乎不变，所以它们单独过不了方向这一步。

\section{一格要进头条，七步都要过}

头条窗口是线性十站点验证面板上的 episode 42、站点 0、第 10 天、dose \(=1\)。反事实只改了 \(b_0\) 的第 10 个滞后，从 \(1.063\) 放到 \(0.063\)。

\begin{center}
\small
\begin{tabular}{@{}clp{0.62\textwidth}@{}}
\toprule
步 & 问的是 & 这一格 \\
\midrule
1 & 哪些坐标被改了 & 只有 \(b_0\) 的 lag 10。协变量、干扰块、其余滞后保持观测值 \\
2 & 编辑坐标还在训练盒子里吗 & 训练 \(z\) 从 \(+2.53\) 到 \(+0.09\)，阈值是 4 \\
3 & 改之前这一格是否已经在外面 & 异常日的 \(z=+2.53\)，也在盒子里。替换是往训练中心走 \\
4 & 没被改的背景是否都在盒子里 & 整条向量的 \(s_\infty=2.58\)，落在人口统计 \\
5 & 能不能把背景钳回盒子再当头条 & 不能。钳人口统计或干扰块就换了问题。这一格不需要钳 \\
6 & 均值差的方向是否和这次编辑一致 & 观测减反事实，signed \(=+1.009\) \\
7 & 无干预的那一对是不是很小 & \(\Delta_{\mathrm{base}}=0.021\)，\(\Delta^{\mathrm{cf}}=0.674\) \\
\bottomrule
\end{tabular}
\end{center}

写成头条的句子是：参考路径上的购买比观测低大约 1 个指数单位。\(\Delta^{\mathrm{cf}}=0.674\) 和 \(\delta^{2}=1.192\) 是距离，方向只记在 signed 里。

1920 格里有 2 格的 \(s_\infty\) 大于 4，都在第 4 步被拒。拒绝它们的不是这次编辑：观测条件和反事实条件上的越界坐标相同。

\begin{center}
\small
\begin{tabular}{@{}llrrr@{}}
\toprule
窗口 & 越界坐标 & 观测 \(z\) & 反事实 \(z\) & 原始 \(\Delta\) \\
\midrule
episode 47，站点 5，第 6 天 & 人口统计 & 4.14 & 4.14 & --- \\
episode 43，站点 7，第 10 天 & 干扰块 \(b_1\) lag 10 & 4.18 & 4.18 & 0.019 \\
\bottomrule
\end{tabular}
\end{center}

背景出了支持，两批之差会把这一格的编辑和模型在盒子外的误差混在一起。钳制只允许用在编辑坐标自己越过 4 的时候，钳过的数不是头条。盒子阈值是坐标标准化绝对值 4。

\section{五层才进入业务动作}

通过线在跑实验之前写死。真实面板的第一层只记假设检查。

\begin{center}
\small
\begin{tabular}{@{}lp{0.72\textwidth}@{}}
\toprule
层 & 通过线 \\
\midrule
L1 排序 & 至少三种构造过线 \\
L2 时间 & placebo、峰值位置、门控都过 \\
L3 校准 & PIT、Wasserstein、HSIC 都过 \\
L4 代理 & 至少三个代理同向 \\
L5 敏感 & 变异系数和 E-value 都过 \\
\bottomrule
\end{tabular}
\end{center}

观测校准是必要条件。PIT 没过，或者核宽的变异系数过了 0.20，分数就停在低可信度，不进入业务动作。

\section{四张面板停在哪里}

\begin{center}
\small
\begin{tabular}{@{}lcccccc@{}}
\toprule
面板 & L1 & L2 & L3 & L4 & L5 & 可信度 \\
\midrule
长记忆模拟，峰值应在第 10 天 & 通过 & 未过 & 未过 & 通过 & 未过 & 低 \\
短记忆模拟，峰值应在第 1 天 & 通过 & 通过 & 未过 & 通过 & 未过 & 低 \\
北京 PM2.5，2014--2017 & --- & --- & --- & --- & --- & 停止 \\
美国各州疫情，2020-03 至 2021-06 & 假设未支持 & 未过 & 未过 & 不确定 & 未过 & 低 \\
\bottomrule
\end{tabular}
\end{center}

长记忆的归一化 \(\mathrm{MMD}^{2}\) 在假设滞后上的扫描是：第 9 天 \(224\times\)，第 10 天 \(206\times\)，第 15 天仍有 \(10.5\times\)。声明的真值是 10 天，峰值落在第 9 天，placebo 也未过，所以时间这一层未过。短记忆的峰值在第 1 天，\(2.10\times\)，placebo 滞后都低于 2，时间这一层通过。疫情面板整条扫描贴着 1，第 10 天是 \(1.00\times\)，扫描的最大值在第 1 天。

\begin{center}
\small
\begin{tabular}{@{}lrr@{}}
\toprule
滞后 & 长记忆 & 短记忆 \\
\midrule
1 & 0.10 & 2.10 \\
2 & 1.84 & 1.36 \\
3 & 0.22 & 1.37 \\
4 & 2.71 & 1.54 \\
5 & 4.59 & 1.32 \\
6 & 2.63 & 1.34 \\
7 & 45.52 & 1.27 \\
8 & 196.6 & 1.85 \\
9 & 224.1 & 1.48 \\
10 & 205.73 & 1.27 \\
11 & 0.17 & 1.29 \\
12 & 0.31 & 1.29 \\
13 & 0.39 & 1.28 \\
14 & 1.68 & 1.27 \\
15 & 10.54 & 1.27 \\
\bottomrule
\end{tabular}
\end{center}

种进去的相对排序，长记忆上三组构造通过：剂量反应的一致率 0.90、Spearman 0.90；重叠度 0.83、0.80；外生隔离的效应量 \(-2.51\)，\(p=0.034\)。单调的「越近越大」和 lag-10 峰值不相容，Spearman 0.16，这一组记在假设上。短记忆的时间距离 Spearman 是 \(-0.685\)，差 \(0.015\) 没到 \(-0.7\)，这一组记失败，另外三组通过，第一层仍然通过。

\begin{center}
\small
\begin{tabular}{@{}lrr@{}}
\toprule
校准与敏感 & 长记忆 & 短记忆 \\
\midrule
Wasserstein / 上限 & 0.082 / 0.095 & 0.078 / 0.093 \\
HSIC 比值 & 0.87 & 0.86 \\
PIT & 四维都拒绝均匀 & 第 3 维 \(p=5\times 10^{-5}\) \\
核宽变异系数，线是 0.20 & 0.71 & 0.25 \\
架构变异系数，线是 0.30 & 0.23 & 0.00 \\
E-value，线是 2 & 19.9 & 32.9 \\
\bottomrule
\end{tabular}
\end{center}

北京，东四看古城，滞后 10 天。五段异常的排序一致率 0.30，Spearman \(-0.60\)，主查询倍数 0.30。协议在这里停止，没有训练扩散模型。10 天前的古城 PM2.5 不能按异常强弱排东四当天的小时分布。

美国疫情，南部各州的对数发病率对东北部 10 天前的区域发病率。归一化倍数 1.00，区间盖在一个没有超过基线的原始 MMD 上。业务窗口里南部确实异常（\(z=2.49\)）。这条 10 天的跨区反事实没有解释它。观测 PIT 和 Wasserstein 都失败，这个「没有超过基线」本身也不可信。

\section{这套逻辑现在能说什么}

模拟面板说明协议能抓住种进去的剂量、重叠和隔离，也能在短记忆上把峰值放回第 1 天。持续性会把长滞后的峰值从第 10 天挪到第 9 天。扩散模型的 PIT 还没均匀。MMD 对核宽的变异系数超过 0.20。这三件事没过之前，分数不进入业务动作。

购买行为的归因因此落在一个窗口、一个坐标上：哪一格被改了，编辑之后还在不在训练盒子里，两批均值差的符号是否和这次编辑一致，无干预的那一对是否小到可以忽略。多个窗口的 \(\rho\) 不合成一个总量。北京那一格没有模型，就不写句子。

\end{document}
